\documentclass{article}
\usepackage[margin=2cm]{geometry}
\usepackage{amsmath,amsthm,graphicx,enumitem,tikz}
\newtheorem{theorem}{Theorem}
\newcommand{\project}{TejTeX}
\begin{document}
\section{Introduction}\label{sec:intro}
This is \project. See Section~\ref{sec:method}, Figure~\ref{fig:a}, Table~\ref{tab:a} and Equation~\eqref{eq:a}.\footnote{A footnote.}
Prior work exists~\cite{knuth84}.

\section{Method}\label{sec:method}
\subsection{Lists}
\begin{enumerate}
  \item First
  \begin{enumerate}
    \item Nested
  \end{enumerate}
  \item Second
\end{enumerate}
\begin{enumerate}[label=(\roman*)]
  \item Roman one
  \item Roman two
\end{enumerate}
\begin{itemize}
  \item Bullet
\end{itemize}

\vspace{1cm}

\begin{figure}
  \centering
  \includegraphics[width=0.5\textwidth]{figure.png}
  \caption{A figure}\label{fig:a}
\end{figure}

\begin{table}
  \centering
  \caption{A table}\label{tab:a}
  \begin{tabular}{ll}
    A & B \\
    1 & 2 \\
  \end{tabular}
\end{table}

\begin{equation}\label{eq:a}
  E = mc^2
\end{equation}

\begin{theorem}\label{thm:a}
  Every test passes.
\end{theorem}

\begin{tikzpicture}
  \draw (0,0) -- (1,1);
\end{tikzpicture}

\begin{thebibliography}{9}
\bibitem{knuth84} D. Knuth. The TeXbook. 1984.
\end{thebibliography}
\end{document}
